% F -- The eight published works this report leans on, grouped by what each gave.
\begin{figure}[!htbp]
\centering
\fitwidth{%
\begin{tikzpicture}[font=\small,
  band/.style={anchor=north west,align=left,text width=2.35cm,
               font=\footnotesize\bfseries,text=clmodel},
  card/.style={anchor=north west,draw=clrule,fill=white,rounded corners=3pt,
               align=left,inner sep=5pt,text width=3.42cm,minimum height=3.15cm,
               font=\scriptsize\linespread{0.95}\selectfont},
]
\foreach \yy in {0,-3.45,-6.90,-10.35}{ \draw[clrule!60] (0,\yy) -- (14.8,\yy); }

% ---------------- band 1 -------------------------------------------------
\node[band] at (0,-0.14) {How much room\\ there is};
\node[card] at (2.70,-0.14)
  {{\bfseries\color{clmodel}Song et al., 2010}\,\cite{song2010limits}\\[3pt]
   45{,}000 people, three months. The best score any predictor could ever
   reach: \textbf{93\,\%}.\\[3pt]
   {\itshape\color{clgain}Gave me the ceiling, and the hard hours: midday and
   early evening.}};
\node[card] at (6.80,-0.14)
  {{\bfseries\color{clmodel}Ikanovic \& Mollgaard, 2017}\,\cite{ikanovic2017alternative}\\[3pt]
   Delete the moments where nobody moved, measure again: the ceiling falls to
   \textbf{71\,\%}.\\[3pt]
   {\itshape\color{clgain}Gave me my competitor.}};

% ---------------- band 2 -------------------------------------------------
\node[band] at (0,-3.59) {What to\\ predict with};
\node[card] at (2.70,-3.59)
  {{\bfseries\color{clmodel}Luca et al., 2021}\,\cite{luca2021survey}\\[3pt]
   The survey of the field. The time of day is a standard input; three depths
   is the standard report.\\[3pt]
   {\itshape\color{clgain}Gave me the input, and the way to report.}};
\node[card] at (6.80,-3.59)
  {{\bfseries\color{clmodel}Wang et al., 2023}\,\cite{wang2023llmmob}\\[3pt]
   A person's recent stops written out as English and handed to a large model.
   Nothing trained.\\[3pt]
   {\itshape\color{clgain}The branch I did not take.}};
\node[card] at (10.90,-3.59)
  {{\bfseries\color{clmodel}Qin et al., 2025}\,\cite{qin2025mobllm}\\[3pt]
   A small open model retrained the same cheap way, over six mobility
   datasets.\\[3pt]
   {\itshape\color{clgain}The published work mine is closest to.}};

% ---------------- band 3 -------------------------------------------------
\node[band] at (0,-7.04) {Who to\\ predict for};
\node[card] at (2.70,-7.04)
  {{\bfseries\color{clmodel}Pappalardo et al., 2015}\,\cite{pappalardo2015returners}\\[3pt]
   A population is not one scale of mobility. It splits into \emph{returners}
   and \emph{explorers}.\\[3pt]
   {\itshape\color{clgain}Same idea; my axis is \emph{when} they move, not how
   far.}};
\node[card] at (6.80,-7.04)
  {{\bfseries\color{clmodel}Jeong et al., 2016}\,\cite{jeong2016cluster}\\[3pt]
   One person's history is too short to learn from. Group people who resemble
   each other instead.\\[3pt]
   {\itshape\color{clgain}The mechanism behind Chapter~\ref{ch:groups}.}};

% ---------------- band 4 -------------------------------------------------
\node[band] at (0,-10.49) {What exactly\\ is predicted};
\node[card] at (2.70,-10.49)
  {{\bfseries\color{clmodel}Chen et al., 2013}\,\cite{chen2013nextcell}\\[3pt]
   Predicts the next \emph{antenna} rather than the next place, from radio
   measurements.\\[3pt]
   {\itshape\color{clgain}My target, without their radio measurements.}};
\node[anchor=north west,draw=clbase,dashed,fill=clbase!8,rounded corners=3pt,
      align=left,inner sep=6pt,text width=7.4cm,minimum height=3.15cm,
      font=\footnotesize] at (6.80,-10.49)
  {\textbf{The gap this report sits in.} Everything above predicts a
   \emph{place}, except the last card, which predicts an antenna but has radio
   measurements to do it with. Predicting an antenna from a trace of antennas
   alone inherits a drawback from each side, and Chapter~\ref{ch:data} measures
   the price.};
\end{tikzpicture}}
\caption{The eight published works, and what each one gave me}
\label{fig:sota-map}
\figtext{%
\figpara{Taxonomy of state-of-the-art mobility prediction paradigms in Figure~\ref{fig:sota-map}.}{Top to bottom in Figure~\ref{fig:sota-map}, from
the frame every later number is read in, down to the exact object being predicted.
The first band fixes how much room a predictor has at all, and is why this report
is written against a one-line rule; the second places the instrument; the third is
the ancestry of the only substantial gain obtained; the fourth marks the
boundary between the published task and the one actually evaluated. The italic
line on each card in Figure~\ref{fig:sota-map} says what that work gave to this project.}}

\end{figure}
